\section{Conditional arithmetic routes and exponent ledgers}
\label{sec:tpc68-conditional-arithmetic}

This section states transparent sufficient premises for an asymptotic
family \(X\to\infty\).  None of the premises below is proved in this paper.
Every \(o(1)\) is required to be uniform over all active rows, row pairs,
blocks, sides, and components in the declared physical scope.

\subsection{Route P: uniform literal pair defects}

\begin{proposition}[Conditional pair route]
\label{prop:tpc68-conditional-pair}
Assume that the literal pair-defect matrix satisfies
\begin{align}
 \Delta_P&\le X^{\delta+o(1)},
 \label{eq:tpc68-L2-pair-degree}\\
 P_{mn}&\le X^{-\tau+o(1)}
 \qquad(m\ne n).
 \label{eq:tpc68-L2-pair-saving}
\end{align}
Then
\begin{equation}
 \|H\|
 \le X^{\delta-\tau+o(1)}.
 \label{eq:tpc68-pair-exponent-output}
\end{equation}
Thus this route gives a subpower certificate only when
\(\tau\ge\delta\), and gives decay only when \(\tau>\delta\).
\end{proposition}

\begin{proof}
Apply \cref{cor:tpc68-pair-degree} with the two assumed exponents.
\end{proof}

The hypothesis concerns the completed fixed-\(h_0\) pair entries with
their literal complex weights.  An averaged shift estimate, a density-one
set of row pairs, or a reduced-cell support count does not imply
\eqref{eq:tpc68-L2-pair-saving}.

\subsection{Route S: uniform star cancellation}

\begin{proposition}[Conditional star route]
\label{prop:tpc68-conditional-star}
Assume \eqref{eq:tpc68-L2-pair-degree} and
\begin{equation}
 \max_m\left(\sum_{n\ne m}P_{mn}^2\right)^{1/2}
 \le X^{-\sigma+o(1)}.
 \label{eq:tpc68-L2-star-saving}
\end{equation}
Then
\begin{equation}
 \|H\|
 \le X^{\delta/2-\sigma+o(1)}.
 \label{eq:tpc68-star-exponent-output}
\end{equation}
This route gives a subpower certificate when
\(\sigma\ge\delta/2\), and decay when \(\sigma>\delta/2\).
\end{proposition}

\begin{proof}
Use \cref{cor:tpc68-pair-to-star}.
\end{proof}

An estimate for one scalar sum \(\sum_nH_{mn}\) is insufficient by
\cref{prop:tpc68-local-star-extremizer}.  The premise must control the
entire \(\ell^2\)-star for every actual row.

\subsection{Route W: complete operator walks}

\begin{proposition}[Conditional operator-walk route]
\label{prop:tpc68-conditional-walk}
Let \(q_X\ge1\) be integers.  If
\begin{equation}
 \Omega_{2q_X}(H)
 \le
 X^{2q_X\nu+o(q_X\log X)},
 \label{eq:tpc68-L2-walk-hypothesis}
\end{equation}
then
\begin{equation}
 \|H\|\le X^{\nu+o(1)}.
 \label{eq:tpc68-walk-exponent-output}
\end{equation}
The same conclusion follows from the analogous bound
\[
 \tr H^{2q_X}
 \le X^{2q_X\nu+o(q_X\log X)}.
\]
\end{proposition}

\begin{proof}
Take \(2q_X\)-th roots in
\cref{thm:tpc68-even-walk-certificate,thm:tpc68-even-moment}.
\end{proof}

When \(q_X\to\infty\), the error term in
\eqref{eq:tpc68-L2-walk-hypothesis} must be uniform at that growing walk
length.  A heuristic square-root cancellation for independent signs, a
fixed-\(q\) estimate with uncontrolled constants, or a count of unsigned
support walks is not the stated premise.

\subsection{Missing mass, represented scale, and strict reassembly}

Suppose the physical represented and missing maps are now placed in the
same declared scope.  Let \(\mathcal E_X(z)\) be one fixed reference
energy, put \(S_A=(\Ker A)^\perp\), and assume uniformly for
\(z\in S_A\) that
\begin{align}
 \langle D_Mz,z\rangle
 &\le X^{-\sigma_M+o(1)}\mathcal E_X(z),
 \label{eq:tpc68-L2-missing-mass}\\
 \|Az\|^2
 &\ge X^{-\lambda_{\rm rep}+o(1)}\mathcal E_X(z),
 \label{eq:tpc68-L2-represented-lower}\\
 \kappa_M
 &\le X^{\nu_M+o(1)}.
 \label{eq:tpc68-L2-kappa-exponent}
\end{align}
The certificate lattice may be used to prove the last inequality only if
one of the preceding L2 routes is actually established.

\begin{proposition}[Conditional strict exponent gate]
\label{prop:tpc68-strict-exponent-gate}
Under
\eqref{eq:tpc68-L2-missing-mass}--\eqref{eq:tpc68-L2-kappa-exponent},
\begin{equation}
 \varepsilon_X^2
 :=
 \sup_{z\in S_A\setminus\{0\}}
 \frac{\|Mz\|^2}{\|Az\|^2}
 \le
 X^{\nu_M-\sigma_M+\lambda_{\rm rep}+o(1)}.
 \label{eq:tpc68-relative-exponent}
\end{equation}
Hence
\begin{equation}
 \boxed{\sigma_M>\nu_M+\lambda_{\rm rep}}
 \label{eq:tpc68-strict-reassembly-exponent}
\end{equation}
is sufficient for \(\varepsilon_X=o(1)\).  If, in addition,
\begin{equation}
 \|C_N\|=o(1),
 \label{eq:tpc68-L2-shorting-small}
\end{equation}
then
\[
 \delta(A,M)\ge1-o(1).
\]
\end{proposition}

\begin{proof}
Combine the three assumed inequalities to obtain
\eqref{eq:tpc68-relative-exponent}.  The strict inequality
\eqref{eq:tpc68-strict-reassembly-exponent} makes its right side tend to
zero.  Then \cref{cor:tpc68-canonical-perturbation} and
\eqref{eq:tpc68-L2-shorting-small} give the final assertion.
\end{proof}

Equality in \eqref{eq:tpc68-strict-reassembly-exponent} gives only a
subpower upper bound and does not imply that the perturbation tends to
zero.  Also, no estimate for \(C_N\) follows formally from
\(\kappa_M\): the nuisance output \(B(N)\) is a separate geometric object.

\subsection{Endpoint ledger}

The exact L0 identities and the L1 carrier crosswalk have zero exponent
cost.  Every positive exponent in
\(\delta,\tau,\sigma,\nu_M,\sigma_M,\lambda_{\rm rep}\), every block
summation, and every comparison between reference energies must be charged
once in the common endpoint ledger.  If the literal physical family forces
total loss at least \(1/400\), the endpoint route must stop.  Reallocating
the same loss among pair, star, walk, or reassembly notation does not repair
the budget.
